% \usepackage[utf8]{inputenc}
% \usepackage[natbibapa, maxcitenames=1]{apacite}
% \APACsetEtAl{1}
% \bibliographystyle{apacite}

% \usepackage[utf8]{inputenc}
\usepackage[T1]{fontenc}
\usepackage[spanish, es-tabla]{babel} % reemplaza "cuadro" por "tabla"
% Otros paquetes que necesites cargar

\usepackage[style=apa, backend=biber]{biblatex}
\addbibresource{3_3_BIBLIOGRAFIA/library.bib}
 % Asegúrate de que 'library.bib' sea el nombre correcto de tu archivo .bib


 \usepackage{csquotes}
% \setlength{\bibsep}{5pt plus 0.3ex} %Espaciamiento en la bibliografía
\setcounter{secnumdepth}{3} % Numera los subsubsecciones
\usepackage[T1]{fontenc}
\usepackage[spanish, es-tabla]{babel} %% reemplaza "cuadro" por "tabla"
\decimalpoint  %Cambia coma por punto
% \usepackage{mathptmx}
\usepackage{layouts} %Saber ancho de hoja.
% \usepackage{amssymb}
\usepackage{graphicx}
\usepackage{changepage} %Agregar identación
\usepackage{newtxtext}
% \usepackage{newtxmath}
\usepackage{setspace}
\usepackage[margin=1in]{geometry} % Margenes
\usepackage{float} % Imágenes con [H]
\usepackage{tabularx} % para tablas de ancho ajustado
\usepackage{fancyhdr} %Encabezado y Pie de página (1)
\pagestyle{fancy} %Encabezado y Pie de página (2)
\usepackage{lipsum} %Crear texto RAMDOM
\renewcommand{\labelitemi}{$\bullet$} %Circulos para viñetas
\usepackage{titlesec} %Titulos de SECCIONES
\usepackage{tocloft} %Titulos de ÍNDICES
\usepackage[colorlinks=true,linkcolor=negro,citecolor = negro]{hyperref}
\usepackage[mathbf=sym]{unicode-math} % Mantener fuentes matemáticas
\makeatletter  %Comando para REDUCIR ERRORES
\usepackage{multirow} % Agregar TABLAS 
\usepackage{array} % Dar formato a las TABLAS
\usepackage{subcaption} % Insertar SubImagenes
\usepackage{tikz} % Diagrama de Flujo
\usetikzlibrary{calc,positioning,shapes.geometric,shapes.symbols,shapes.misc}
\usepackage{chngcntr}
\counterwithout{figure}{chapter}
% Insertar formas para diagramas de flujo
\tikzstyle{startstop} = [rectangle, rounded corners, minimum width=3cm, minimum height=0.5cm,text centered, draw=black]
\tikzstyle{io} = [trapezium, trapezium left angle=70, trapezium right angle=110, minimum width=3cm, minimum height=0.5cm, text centered, text width=3cm, draw=black]
\tikzstyle{process} = [rectangle, minimum width=3cm, minimum height=0.5cm, text centered, text width=4cm, draw=black]
\tikzstyle{decision} = [diamond, minimum width=3cm, minimum height=0.5cm, text centered, draw=black]
\tikzstyle{loop} = [chamfered rectangle,chamfered rectangle xsep=2cm,draw=black]
\tikzstyle{arrow} = [->,>=stealth]
\tikzstyle{line}=[draw]
\usepackage{listings}
\usepackage{color}
\usepackage{ragged2e}
\usepackage{booktabs}
\usepackage{xstring}
%New colors defined below
\definecolor{codegreen}{rgb}{0,0.6,0}
\definecolor{codegray}{rgb}{0.5,0.5,0.5}
\definecolor{codepurple}{rgb}{0.2,0,1}
\definecolor{codeRojo}{rgb}{0.7,0,0.3}
\definecolor{backcolour}{rgb}{1.0, 1.0, 1.0}

%Code listing style named "mystyle"
\lstdefinestyle{mystyle}{
  backgroundcolor=\color{backcolour},  
  commentstyle=\color{codegreen},
  keywordstyle=\color{codeRojo},
  numberstyle=\tiny\color{codegray},
  stringstyle=\color{codepurple},
  basicstyle=\footnotesize,
  breakatwhitespace=false,         
  breaklines=true,                 
  captionpos=b,                    
  keepspaces=true,                 
  numbers=left,                    
  numbersep=5pt,                  
  showspaces=false,                
  showstringspaces=false,
  showtabs=false,                  
  tabsize=2
}

%"mystyle" code listing set
\lstset{style=mystyle}
\newenvironment{MyFont}{\fontfamily{ugm}\selectfont}{\par}
\usepackage{changepage} %Agregar espacio a Listing

% Centrado del título del ÍNDICE / LISTA DE FIGURAS / LISTA DE CUADROS
\renewcommand{\cfttoctitlefont}{\hfill \normalfont\normalsize\bfseries}
\renewcommand{\cftaftertoctitle}{\hfill}
\renewcommand{\cftlottitlefont}{\hfill\normalfont\normalsize\bfseries}
\renewcommand{\cftafterlottitle}{\hfill}
\renewcommand{\cftloftitlefont}{\hfill\normalfont\normalsize\bfseries}
\renewcommand{\cftafterloftitle}{\hfill}




% --------------------------
% Nuevo Contador para Capítulos Especiales
% --------------------------
\newcounter{ChapterRoman}
\renewcommand{\theChapterRoman}{\Roman{ChapterRoman}}

% --------------------------
% Nuevo Comando para Capítulos Especiales
% --------------------------
% \newcommand{\Chapter}[1]{
%     \refstepcounter{ChapterRoman} % Incrementa el contador
%     \chapter*{CAPÍTULO \theChapterRoman\ \ #1} % Imprime el título del capítulo con el prefijo y el contador
%     \addcontentsline{toc}{chapter}{CAPÍTULO \theChapterRoman: #1} % Agrega la entrada al Índice
% }


\newcommand{\Chapter}[1]{
    \clearpage % Asegura que empezamos en una página nueva
    \refstepcounter{ChapterRoman} % Incrementa el contador

    \thispagestyle{empty} % Elimina encabezados y pie de página para esta página
    \null\vfill % Centra el contenido verticalmente
    \begin{center} % Centra el contenido horizontalmente
        \Large\bfseries CAPÍTULO \theChapterRoman\ \\[1em] #1 % Título del capítulo
    \end{center}
    \vfill\null % Centra el contenido verticalmente
    \addcontentsline{toc}{chapter}{CAPÍTULO \theChapterRoman: #1} % Agrega la entrada al Índice
    \clearpage % Comienza una nueva página después del título del capítulo
}

% ----------------------------------
% FORMATEO DE ENCABEZADOS DE SECCIÓN
% ----------------------------------

% Nivel 1 (\chapter)
% - Formato: fuente de tamaño normal, negrita, centrado
% - El número del capítulo seguido por un punto y un espacio de 0.5em
% - El título del capítulo en mayúsculas
\titleformat{\chapter}[block]{\normalfont\normalsize\bfseries\centering}{\thechapter.}{0.5em}{\MakeUppercase}
\titlespacing*{\chapter}{0pt}{12pt}{5pt} % Espaciado: Antes del título 12pt, después del título 5pt

% Nivel 2 (\section)
% - Formato: fuente de tamaño normal, negrita, alineada a la izquierda
% - El número de la sección seguido por un espacio de 0.5em
\titleformat{\section}[block]{\normalfont\normalsize\bfseries\raggedright}{\thesection}{0.5em}{}
\titlespacing*{\section}{0pt}{12pt}{5pt} % Espaciado: Antes del título 12pt, después del título 5pt

% Nivel 3 (\subsection)
% - Formato: fuente de tamaño normal, negrita, cursiva, alineada a la izquierda
% - El número de la subsección seguido por un espacio de 0.5em
\titleformat{\subsection}[block]{\normalfont\normalsize\bfseries\itshape\raggedright}{\thesubsection}{0.5em}{}
\titlespacing*{\subsection}{0pt}{12pt}{5pt} % Espaciado: Antes del título 12pt, después del título 5pt

% Nivel 4 (\subsubsection)
% - Formato: fuente de tamaño normal, negrita, alineada a la izquierda, estilo "runin" (seguirá con el texto en la misma línea)
% - El número de la subsubsección seguido por un espacio de 0.5em
% - Termina con un punto y dos espacios
\titleformat{\subsubsection}[runin]{\normalfont\normalsize\bfseries\raggedright}{\thesubsubsection}{0.5em}{}[. \quad]
\titlespacing*{\subsubsection}{1.27cm}{12pt}{0pt} % Espaciado: Antes del título 1.27cm, después del título 0pt

% Nivel 5 (\paragraph)
% - Formato: fuente de tamaño normal, negrita, cursiva, alineada a la izquierda, estilo "runin" (seguirá con el texto en la misma línea)
% - El número del párrafo seguido por un espacio de 0.5em
% - Termina con un punto y dos espacios
\titleformat{\paragraph}[runin]{\normalfont\normalsize\bfseries\itshape\raggedright}{\theparagraph}{0.5em}{}[. \quad]
\titlespacing*{\paragraph}{1.27cm}{12pt}{0pt} % Espaciado: Antes del título 1.27cm, después del título 0pt

% ------------------------------------
% CAMBIOS EN LA NUMERACIÓN DE CAPÍTULOS
% ------------------------------------

% Cambia el formato de numeración de capítulos a números árabes
% (descomenta la línea siguiente para usar números romanos)
% \renewcommand{\thechapter}{\Roman{chapter}}
\renewcommand{\thechapter}{\arabic{chapter}}

% Cambia el formato de numeración de ecuaciones para incluir el número de capítulo
\renewcommand{\theequation}{\arabic{chapter}.\arabic{equation}} 

% Cambia el formato de numeración de secciones para incluir el número de capítulo
\renewcommand{\thesection}{\arabic{chapter}.\arabic{section}}  


% Espaciado entre PÁRRAFOS y SANGRÍA
% \setlength{\parskip}{5 pt}
\setlength{\parindent}{1.27cm}
\doublespacing




% Cambiar titulo de bibliografía
\addto\captionsspanish{\renewcommand{\bibname}{\centering Referencias bibliográficas}}


% Posicionamiento vertical de TOC, LOT and LOF

\setlength{\cftbeforelottitleskip}{1pt} 
\renewcommand{\cftafterlottitleskip}{12pt}

\setlength{\cftbeforeloftitleskip}{1pt} 
\renewcommand{\cftafterloftitleskip}{12pt}

\setlength{\cftbeforetoctitleskip}{16pt} 
\renewcommand{\cftaftertoctitleskip}{12 pt}

%-------------------------------------
% CONFIGURACIÓN DE FIGURAS
%-------------------------------------

% Ajuste de la nomenclatura para las figuras en español:
% Reemplaza el nombre por defecto "Figura" en las etiquetas de las imágenes. 
\addto\captionsspanish{\renewcommand{\figurename}{Figura}}

% Configuración del formato del caption (etiqueta) para figuras:
\captionsetup{
   % Alinea el texto del caption a la izquierda.
   justification=raggedright,
   % La justificación se aplica incluso si el texto del caption es corto.
   singlelinecheck=false,
   % Añade un punto después del número en el caption (por ejemplo, "Figura 1.").
   labelsep=period,
   % Pone en negrita el texto "Figura Número" en el caption.
   labelfont=bf 
}

% Comando personalizado para insertar figuras:
% Argumentos:
%   #1: Título de la figura.
%   #2: Ruta de la imagen.
%   #3: Descripción o nota asociada a la figura.
\newcommand{\insertfigure}[3]{
    \begin{figure}[H] % El entorno 'figure' con la opción [H] asegura que la figura se coloque exactamente en este punto del documento.
        
        % Inserta el título de la figura, aplicando doble espacio y formato en cursiva.
        \caption{\doublespacing \\ \textit{#1}} 
        
        % Centra e inserta la imagen. La anchura es ajustada al ancho de la línea del texto.
        \centering
        \includegraphics[width=0.7\linewidth]{#2}

        % Añade una nota al pie de la figura, completamente justificada y en cursiva.
        % \doublespacing
        \begin{justify}
            \textit{Nota.} #3
        \end{justify}
        
        % Etiqueta para referencia cruzada. Usa el título de la figura para generar una etiqueta única.
        
        \label{fig:#1}
    \end{figure}
}

% %-------------------------------------
% % CONFIGURACIÓN DE TABLAS
% %-------------------------------------
\counterwithout{table}{chapter}
% Ajuste de la nomenclatura para las tablas en español:
\addto\captionsspanish{\renewcommand{\tablename}{Tabla}}

% Configuración del formato del caption (etiqueta) para tablas:
\captionsetup[table]{
   justification=raggedright,
   singlelinecheck=false,
   labelsep=period,
   labelfont=bf 
}



\newcolumntype{P}[1]{>{\justifying\noindent\arraybackslash}p{#1}}

% \newcommand{\inserttable}[8]{
%     \begin{table}[H]
%         \caption{\doublespacing \\ \textit{#1}}
%         \begin{spacing}{8}
%             \fontsize{#2}{1.5*#2}\selectfont  
%             \begin{tabularx}{\linewidth}{#3} % *{6}{c}P{3.5cm}
%                 \toprule
%                 #4
%                 \bottomrule
%             \end{tabularx}
        
%         \end{spacing}
%         \vspace{1\baselineskip}
%         \textit{Nota.} #5
%         \label{#1}
%     \end{table}
% }


          




% Define que las tablas sean numeradas con números arábigos
\renewcommand{\thetable}{\arabic{table}}  



% Espaciamiento dentro del índice

\setlength{\cftbeforechapskip}{2mm}
\renewcommand\cftchapafterpnum{\vskip6pt}
\renewcommand\cftsecafterpnum{\vskip5pt}
\renewcommand\cftsubsecafterpnum{\vskip5pt}


%Agregar la palabra CAPITULO al TOC, FIGURA a LOF y TABLA al LOT

\renewcommand{\cftchappresnum}{CAPÍTULO }
\renewcommand{\cftchapaftersnum}{:}
\renewcommand{\cftchapnumwidth}{7em}

\renewcommand{\cftfigpresnum}{Figura N° }
\renewcommand{\cftfigaftersnum}{.}
\renewcommand{\cftfignumwidth}{6.85 em}

\renewcommand{\cfttabpresnum}{Tabla N° }
%\renewcommand{\cftfigaftersnum}{:}
\renewcommand{\cfttabnumwidth}{6.5 em}


%Definición de COLORES

\usepackage{xcolor}
\definecolor{granate}{RGB}{113,22,16}
\definecolor{gris}{RGB}{154,153,157}
\definecolor{arena}{RGB}{230,217,170}
\definecolor{azul}{rgb}{0.03, 0.15, 0.4}
\definecolor{negro}{rgb}{0, 0, 0}



%Escribir los INFORMACIÓN PERSONAL Y DEL TRABAJO

%Autor para PIE DE PÁGINA (Respetar mayusculas y minisculas)
\author{Nombre1 Nombre2 Apellido1 Apellido2}

%Autor para CARÁTULA (Siempre en mayuscula y sin saltos de linea)
\authorcaratula{NOMBRE1 NOMBRE2 APELLIDO1 APELLIDO2}

%Título en para el PIE DE PÁGINA (Agregar salto de línea de ser necesario)
\title{Xxxxxxxxxxx Título de la Tesis Xxxxxxxxxxx}

%Título para CARÁTULA (Siempre en mayuscula y sin saltos de linea)
\titlecaratula{XXXXXXXXXXXX TÍTULO DE LA TESIS XXXXXXXXXXXXXX}

%Nombre de la FACULTAD (Siempre en mayuscula)
\facultad{FACULTAD DE INGENIERÍA CIVIL}

%Para obtener el título profesional de ... (Siempre en mayuscula)
\grado{INGENIERO CIVIL}

%Asesor para CARÁTULA (Siempre en mayuscula y sin saltos de linea)
\asesor{DR. NOMBRE DEL ASESOR}

%AÑO para CARÁTULA
\yyearr{2024}

\lstset{
  language=Python,
  basicstyle=\ttfamily\small,
  commentstyle=\color{gray},
  keywordstyle=\color{blue},
  stringstyle=\color{red},
  numberstyle=\tiny\color{gray},
  numbers=left,
  stepnumber=1,
  showstringspaces=false,
  tabsize=4,
  breaklines=true,
  frame=single,
  literate={á}{{\'a}}1 {é}{{\'e}}1 {í}{{\'i}}1 {ó}{{\'o}}1 {ú}{{\'u}}1
           {Á}{{\'A}}1 {É}{{\'E}}1 {Í}{{\'I}}1 {Ó}{{\'O}}1 {Ú}{{\'U}}1
           {ñ}{{\~n}}1 {Ñ}{{\~N}}1 {ü}{{\"u}}1 {Ü}{{\"U}}1
}